\documentclass[twoside]{book}
\topmargin=-30pt
\oddsidemargin=-20pt
\evensidemargin=-20pt
\textwidth=17cm
\textheight=25cm

\usepackage{czech,amsmath,amsfonts,graphics}
%\usepackage{pspicture,picins}
\DeclareMathOperator{\tg}{tg}
\DeclareMathOperator{\arctg}{arctg}
\DeclareMathOperator{\cotg}{cotg}
\DeclareMathOperator{\arccotg}{arccotg}
\DeclareMathOperator{\hod}{hod}
\DeclareMathOperator{\sgn}{sgn}
\makeindex
\begin{document}

\pagestyle{empty}

\def\dr #1{\par\noindent $\bullet$ {\bf #1}}


\dr{Friedmann�v test}
\vspace*{2mm}
\begin{multline*}\begin{array}{l}
X_{1,1},\dots,X_{1,J}  \\
X_{2,1},\dots,X_{2,J}   \\
\hdotsfor{1}  \\
X_{I,1},\dots,X_{I,J}
\end{array}   \end{multline*}
\hspace*{80mm} nez�visl�
\newline $H_0$: $X_{1l}, X_{2l}, \dots, X_{Il}$ maj� stejn� rozd�len�
\newline
postup:
\newline - v ka�d�m bloku (sloupci) p�i�ad�me po�ad� $R_{k,1}, \dots R_{k,J}$ hodnot�m $X_{k,1}, \dots X_{k,J}$
\newline - $ Q=\frac{12}{JI(I+1)} \sum_{l=1}^{I} \left(\sum_{k=1}^{J} R_{kl}\right)^2  -3J(I+1)\approx \chi^2_{I-1}$
\newline $H_0$  zam�tneme $\quad \Longleftrightarrow\quad |Q| \geq k$
\newline \phantom{$H_0$  zam�tneme $\quad \Longleftrightarrow$} $\quad |Q| \geq \chi^2_{I-1}(1-\alpha)$



\end{document}


